
\documentclass{article}
\usepackage{colortbl}
\usepackage{makecell}
\usepackage{multirow}
\usepackage{supertabular}

\begin{document}

\newcounter{utterance}

\centering \large Interaction Transcript for game `ta\_blackjack', experiment `Blackjack{-}v0', episode 1 with BSU{-}SLIM\_\_Qwen3.5{-}9B{-}Base\_\_playpen{-}prm{-}9b{-}best\_of\_n.
\vspace{24pt}

{ \footnotesize  \setcounter{utterance}{1}
\setlength{\tabcolsep}{0pt}
\begin{supertabular}{c@{$\;$}|p{.15\linewidth}@{}p{.15\linewidth}p{.15\linewidth}p{.15\linewidth}p{.15\linewidth}p{.15\linewidth}}
    \# & \multicolumn{2}{c}{Player} && \multicolumn{2}{c}{Game Master} \\
    \hline

    \theutterance \stepcounter{utterance}  
    & & & \multicolumn{4}{p{0.6\linewidth}}{
        \cellcolor[rgb]{0.9,0.9,0.9}{
            \makecell[{{p{\linewidth}}}]{
                \texttt{\tiny{[P1$\langle$GM]}}
                \texttt{You are a competitive game player. Make sure you read the game instructions carefully, and always follow the required format.} \\
\\ 
\texttt{{[}GAME{]} You are playing Blackjack against the dealer.} \\
\texttt{Your goal is to get as close to 21 as possible without going over.} \\
\texttt{On your turn, choose '{[}Hit{]}' to draw another card or '{[}Stand{]}' to hold.} \\
\texttt{J/Q/K = 10 points; A = 11 or 1, whichever is better.} \\
\\ 
\texttt{{[}GAME{]} Hand 1/5} \\
\texttt{Your hand: 5♣, 2♥ (Score: 7)} \\
\texttt{Dealer shows: 4♣} \\
            }
        }
    }
     \\ \\

    \theutterance \stepcounter{utterance}  
    & \multicolumn{4}{p{0.6\linewidth}}{
        \cellcolor[rgb]{0.561,0.737,0.561}{
            \makecell[{{p{\linewidth}}}]{
                \texttt{\tiny{[P1$\rangle$GM]}}
                \texttt{I'll take **Hit**.} \\
\\ 
\texttt{With a current score of 7 and the dealer showing a 4, I have a strong position to improve. The dealer's 4 is a weak start, meaning they have a good chance of busting if they draw multiple cards. Taking another card gives me the best opportunity to reach a higher total (ideally 17 or higher) before the dealer can complete their hand.} \\
\\ 
\texttt{Given I need to get closer to 21 without going over, hitting is the optimal move here.} \\
            }
        }
    }
    & & \\ \\

    \theutterance \stepcounter{utterance}  
    & & & \multicolumn{4}{p{0.6\linewidth}}{
        \cellcolor[rgb]{0.9,0.9,0.9}{
            \makecell[{{p{\linewidth}}}]{
                \texttt{\tiny{[P1$\langle$GM]}}
                \texttt{{[}GAME{]} You attempted an invalid move. Reason: Invalid action. Use '{[}Hit{]}' or '{[}Stand{]}'. Please resubmit a valid move and remember to follow the game rules to avoid penalties.} \\
            }
        }
    }
     \\ \\

    \theutterance \stepcounter{utterance}  
    & \multicolumn{4}{p{0.6\linewidth}}{
        \cellcolor[rgb]{0.561,0.737,0.561}{
            \makecell[{{p{\linewidth}}}]{
                \texttt{\tiny{[P1$\rangle$GM]}}
                \texttt{{[}Stand{]}} \\
\\ 
\texttt{I apologize for the previous error. With a hard 7 and the dealer showing a 4, I will stand. Standing allows me to see the dealer's hidden card and potential bust, which is a strong strategic position when holding a low score against a weak dealer upcard, and it is the only valid action format required.} \\
            }
        }
    }
    & & \\ \\

    \theutterance \stepcounter{utterance}  
    & & & \multicolumn{4}{p{0.6\linewidth}}{
        \cellcolor[rgb]{0.9,0.9,0.9}{
            \makecell[{{p{\linewidth}}}]{
                \texttt{\tiny{[P1$\langle$GM]}}
                \texttt{{[}GAME{]} Hand 1: you lose. Your final 7, Dealer 17.} \\
\texttt{{[}GAME{]} Hand 2/5} \\
\texttt{Your hand: 8♥, 6♥ (Score: 14)} \\
\texttt{Dealer shows: 8♦} \\
            }
        }
    }
     \\ \\

    \theutterance \stepcounter{utterance}  
    & \multicolumn{4}{p{0.6\linewidth}}{
        \cellcolor[rgb]{0.561,0.737,0.561}{
            \makecell[{{p{\linewidth}}}]{
                \texttt{\tiny{[P1$\rangle$GM]}}
                \texttt{{[}Hit{]}} \\
\\ 
\texttt{My current score is 14, which is a risky position because I'm just one card away from busting (21). The dealer showed an 8, which is a strong card. Statistically, the dealer has a high probability of reaching a final score of 17, 18, 19, 20, or 21 with an 8 upcard. If I stand at 14 and the dealer makes 17 or higher, I lose. Therefore, I must hit to try and improve my hand to a safer total (like 17 or 18) before the dealer reveals their hole card and completes their hand.} \\
            }
        }
    }
    & & \\ \\

    \theutterance \stepcounter{utterance}  
    & & & \multicolumn{4}{p{0.6\linewidth}}{
        \cellcolor[rgb]{0.9,0.9,0.9}{
            \makecell[{{p{\linewidth}}}]{
                \texttt{\tiny{[P1$\langle$GM]}}
                \texttt{{[}GAME{]} Hand 2/5} \\
\texttt{Your hand: 8♥, 6♥, 4♣ (Score: 18)} \\
\texttt{Dealer shows: 8♦} \\
            }
        }
    }
     \\ \\

    \theutterance \stepcounter{utterance}  
    & \multicolumn{4}{p{0.6\linewidth}}{
        \cellcolor[rgb]{0.561,0.737,0.561}{
            \makecell[{{p{\linewidth}}}]{
                \texttt{\tiny{[P1$\rangle$GM]}}
                \texttt{{[}Stand{]}} \\
\\ 
\texttt{My score is now 18, which is a strong position. The dealer showed an 8, meaning they have a high probability of reaching a final score of 17, 18, 19, 20, or 21. Since I already have 18, hitting risks going over 21 (busting) with a single card draw, while standing gives me a chance to win if the dealer ends up with 17 or less, or ties at 18. Standing is the optimal and safe move here.} \\
            }
        }
    }
    & & \\ \\

    \theutterance \stepcounter{utterance}  
    & & & \multicolumn{4}{p{0.6\linewidth}}{
        \cellcolor[rgb]{0.9,0.9,0.9}{
            \makecell[{{p{\linewidth}}}]{
                \texttt{\tiny{[P1$\langle$GM]}}
                \texttt{{[}GAME{]} Hand 2: you win. Your final 18, Dealer 25.} \\
\texttt{{[}GAME{]} Hand 3/5} \\
\texttt{Your hand: 6♦, A♦ (Score: 17)} \\
\texttt{Dealer shows: K♦} \\
            }
        }
    }
     \\ \\

    \theutterance \stepcounter{utterance}  
    & \multicolumn{4}{p{0.6\linewidth}}{
        \cellcolor[rgb]{0.561,0.737,0.561}{
            \makecell[{{p{\linewidth}}}]{
                \texttt{\tiny{[P1$\rangle$GM]}}
                \texttt{{[}Stand{]}} \\
\\ 
\texttt{I have a hard 17, which is a very strong hand. The dealer showing a King (10) is a strong upcard, meaning they have a high probability of reaching 17, 18, 19, 20, or 21. However, hitting on 17 is risky; drawing a 4, 5, 6, 7, 8, 9, or 10 would immediately bust me (going over 21). Even drawing an Ace (counting as 11) would bust me (28), and counting it as 1 would only give me 18, which is less than the dealer's potential best score. Standing is the optimal strategy here to avoid busting while holding a solid total.} \\
            }
        }
    }
    & & \\ \\

    \theutterance \stepcounter{utterance}  
    & & & \multicolumn{4}{p{0.6\linewidth}}{
        \cellcolor[rgb]{0.9,0.9,0.9}{
            \makecell[{{p{\linewidth}}}]{
                \texttt{\tiny{[P1$\langle$GM]}}
                \texttt{{[}GAME{]} Hand 3: you lose. Your final 17, Dealer 18.} \\
\texttt{{[}GAME{]} Hand 4/5} \\
\texttt{Your hand: 7♣, 3♣ (Score: 10)} \\
\texttt{Dealer shows: 6♣} \\
            }
        }
    }
     \\ \\

    \theutterance \stepcounter{utterance}  
    & \multicolumn{4}{p{0.6\linewidth}}{
        \cellcolor[rgb]{0.561,0.737,0.561}{
            \makecell[{{p{\linewidth}}}]{
                \texttt{\tiny{[P1$\rangle$GM]}}
                \texttt{{[}Hit{]}} \\
\\ 
\texttt{I will take another card. With a score of 10 and the dealer showing a 6, standing is too risky. If I stand, the dealer (with a 6 upcard) has a high probability of reaching 17, which would beat my 10. To have a chance of winning, I need to improve my hand to at least 17. Hitting is the necessary move to try and get closer to 21.} \\
            }
        }
    }
    & & \\ \\

    \theutterance \stepcounter{utterance}  
    & & & \multicolumn{4}{p{0.6\linewidth}}{
        \cellcolor[rgb]{0.9,0.9,0.9}{
            \makecell[{{p{\linewidth}}}]{
                \texttt{\tiny{[P1$\langle$GM]}}
                \texttt{{[}GAME{]} Hand 4/5} \\
\texttt{Your hand: 7♣, 3♣, 8♥ (Score: 18)} \\
\texttt{Dealer shows: 6♣} \\
            }
        }
    }
     \\ \\

    \theutterance \stepcounter{utterance}  
    & \multicolumn{4}{p{0.6\linewidth}}{
        \cellcolor[rgb]{0.561,0.737,0.561}{
            \makecell[{{p{\linewidth}}}]{
                \texttt{\tiny{[P1$\rangle$GM]}}
                \texttt{{[}Stand{]}} \\
\\ 
\texttt{My score is 18, which is a very strong hand. The dealer showed a 6, which is a weak upcard compared to high cards like 10 or Ace. Statistically, the dealer has a lower probability of reaching 17 or higher with a 6 upcard than with a 10 or Ace. Since I already have 18, hitting risks busting (going over 21), while standing gives me a high probability of winning if the dealer ends up with 16 or less, or ties at 17/18. Standing is the optimal and safest move here.} \\
            }
        }
    }
    & & \\ \\

    \theutterance \stepcounter{utterance}  
    & & & \multicolumn{4}{p{0.6\linewidth}}{
        \cellcolor[rgb]{0.9,0.9,0.9}{
            \makecell[{{p{\linewidth}}}]{
                \texttt{\tiny{[P1$\langle$GM]}}
                \texttt{{[}GAME{]} Hand 4: you lose. Your final 18, Dealer 20.} \\
\texttt{{[}GAME{]} Hand 5/5} \\
\texttt{Your hand: Q♣, J♦ (Score: 20)} \\
\texttt{Dealer shows: 5♣} \\
            }
        }
    }
     \\ \\

    \theutterance \stepcounter{utterance}  
    & \multicolumn{4}{p{0.6\linewidth}}{
        \cellcolor[rgb]{0.561,0.737,0.561}{
            \makecell[{{p{\linewidth}}}]{
                \texttt{\tiny{[P1$\rangle$GM]}}
                \texttt{{[}Stand{]}} \\
\\ 
\texttt{I have 20, which is an excellent hand and only 1 point away from the natural 21. The dealer is showing a 5, which is a weak upcard. While the dealer has a good chance of eventually reaching 17, 18, 19, or 20, hitting on 20 carries a significant risk of busting immediately if they draw any card worth 2 or more points (which is almost any card except the Ace). Since standing gives me a very high probability of winning (beating a dealer who is likely to settle at 17 or 18) and hitting poses an immediate bust risk, {[}Stand{]} is the optimal move.} \\
            }
        }
    }
    & & \\ \\

    \theutterance \stepcounter{utterance}  
    & & & \multicolumn{2}{p{0.3\linewidth}}{
        \cellcolor[rgb]{0.9,0.9,0.9}{
            \makecell[{{p{\linewidth}}}]{
                \texttt{\tiny{[GM$|$GM]}}
                \texttt{role: Player 0} \\
\texttt{invalid\_move: False} \\
\texttt{turn\_count: 7} \\
\texttt{reason: The game has concluded. Final scores: Dealer: 3, You: 2, Draws: 0} \\
\texttt{reward: 0.4} \\
            }
        }
    }
    & & \\ \\

\end{supertabular}
}

\end{document}
